\paragraph{推前谱的有理交换子刚性、边界协变判据与外自同构层}

\begin{theorem}[推前谱的有理交换子分解与幂等刚性]\label{thm:terminal-windowm-rational-commutant-idempotent-rigidity}
固定 \(m\ge 5\)，令 \(P_m\in M_{|X_m|}(\QQ)\) 为 window-\(m\) 推前核，并设
\[
\chi_{P_m}(\lambda)=(\lambda-1)q_m(\lambda),\qquad q_m\in\QQ[\lambda],
\]
其中 \(q_m\) 在 \(\QQ\) 上不可约。记
\[
K_m:=\QQ[\lambda]/(q_m).
\]
则有 \(\QQ\)-代数同构
\[
\operatorname{Cent}_{M_{|X_m|}(\QQ)}(P_m)\cong \QQ\times K_m.
\]
因此该交换子代数中的幂等元恰为
\[
0,\ I,\ \Pi_m,\ I-\Pi_m,
\]
其中 \(\Pi_m\) 为 \(P_m\) 在特征值 \(1\) 上的谱投影。特别地，除稳态方向与其补空间外，不存在与 \(P_m\) 交换的非平凡有理幂等分解。
若平稳分布满足 \(\pi_m(w)=d^{\mathrm{bin}}_m(w)/2^m\in\QQ\)，则
\[
\Pi_m=\mathbf{1}\,\pi_m^{\mathsf T}\in M_{|X_m|}(\QQ).
\]
\end{theorem}
\leanverified{paper\_terminal\_windowm\_rational\_commutant\_idempotent\_rigidity}
\begin{proof}
令 \(V:=\QQ^{X_m}\)，并把 \(P_m\) 视为 \(V\) 上线性算子。
由 \(\chi_{P_m}=(\lambda-1)q_m\) 且 \(\gcd(\lambda-1,q_m)=1\)，中国剩余分解给出
\[
V=\ker(P_m-I)\oplus\ker\bigl(q_m(P_m)\bigr)=:V_1\oplus V_q.
\]
由于 \((\lambda-1)\) 在特征多项式中只出现一次，\(\dim_{\QQ}V_1=1\)。

另一方面，\(q_m\) 不可约，故 \(\QQ[\lambda]/(q_m)=K_m\) 为域，且 \(V_q\) 自然成为 \(K_m\)-向量空间。
又 \(\dim_{\QQ}V_q=\deg q_m=[K_m:\QQ]\)，从而 \(\dim_{K_m}V_q=1\)。
于是
\[
\End_{\QQ[P_m]}(V_1)\cong \QQ,\qquad
\End_{\QQ[P_m]}(V_q)\cong K_m.
\]
由互素分解下块间同态为零，得到
\[
\operatorname{Cent}_{M_{|X_m|}(\QQ)}(P_m)
=\End_{\QQ[P_m]}(V)
\cong \QQ\times K_m.
\]

域直积中的幂等元只能在每个分量取 \(0\) 或 \(1\)，故恰得四个幂等元，对应于
\(0,I,\Pi_m,I-\Pi_m\)。
这说明与 \(P_m\) 交换的有理精确粗粒化仅可能是平凡分解与稳态/非稳态二分。
若 \(\pi_m(w)=d^{\mathrm{bin}}_m(w)/2^m\)（在 \(m=6\) 的情形由定理 \ref{thm:terminal-foldbin6-pushforward-markov} 明确给出），则 \(\Pi_m=\mathbf{1}\,\pi_m^{\mathsf T}\) 的条目均为有理数。证毕。
\end{proof}

\begin{corollary}[window-\(6\) 的有理可观测压缩平凡性]\label{cor:terminal-window6-rational-observable-compression-triviality}
在 \(V_{\QQ}:=\QQ^{X_6}\) 上，若 \(E\in M_{21}(\QQ)\) 满足
\[
E^2=E,\qquad EP_6=P_6E,
\]
则
\[
E\in\{0,\ I,\ \Pi_6,\ I-\Pi_6\}.
\]
因此，任何把 \(\mathfrak{sl}(V)\) 压缩到低秩 Levi 骨架的非平凡线性读出，不可能仅由与 \(P_6\) 交换的有理可观测量实现；其实现必须引入额外的可观测压缩/协议约束。
\end{corollary}
\begin{proof}
定理 \ref{thm:terminal-window6-pushforward-charpoly-galois} 给出 \(q_6\) 在 \(\QQ\) 上不可约，故可将定理 \ref{thm:terminal-windowm-rational-commutant-idempotent-rigidity} 应用于 \(m=6\)，立即得到幂等元分类。
后半句与命题 \ref{prop:window6-lowrank-needs-observable-compression} 相容：仅靠交换子闭包与有理可观测压缩不能产生低秩 Levi 抽取。证毕。
\end{proof}

\begin{proposition}[符号协变二值标号的存在判据与 boundary 分层分岔]\label{prop:terminal-sign-covariant-label-boundary-bifurcation}
令 \(F\) 为 \(d\) 点集合，\(\mathfrak S_d\) 自然作用于 \(F\)。
存在映射 \(\chi:F\to\{\pm1\}\) 满足
\[
\chi(\sigma\cdot x)=\mathrm{sgn}(\sigma)\chi(x)\qquad(\sigma\in\mathfrak S_d,\ x\in F)
\]
当且仅当 \(d=2\)。

在 boundary uplift 纤维口径下，该判据给出：
\[
m=6:\ \#(\mathrm{Fold}^{\mathrm{bin}}_6)^{-1}(w)=2,\ \Delta_6=\{0,34\}\ \Rightarrow\ \chi\ \text{可实现},
\]
\[
m=7,8:\ \#(\mathrm{Fold}^{\mathrm{bin}}_m)^{-1}(w)=3,\ \Delta_7=\{0,55,89\},\ \Delta_8=\{0,89,144\}\ \Rightarrow\ \chi\ \text{不可实现}.
\]
\end{proposition}
\begin{proof}
首个“当且仅当”结论即推论 \ref{cor:bdry-symmetric-group-sign-twisted-label-d2} 的等价写法。
对 boundary uplift，表 \ref{tab:foldbin_boundary_lift_m6_m8} 给出 \(m=6\) 时每条 boundary 纤维为二层、\(m=7,8\) 时为三层；
并且 \(m=6\) 的位移对 \(\{0,34\}\) 与推论 \ref{cor:terminal-foldbin6-boundary-pure-f9-alias} 一致。
将 \(d=2\) 与 \(d=3\) 代入上述判据即得。证毕。
\end{proof}

\begin{theorem}[window-\(6\) fold-gauge 群的自同构分解与外自同构结构]\label{thm:terminal-window6-fold-gauge-aut-out-structure}
\leanverified{paper\_terminal\_window6\_fold\_gauge\_aut\_out\_structure}
设
\[
\mathcal G^{\mathrm{bin}}_6\cong S_2^{\,8}\times S_3^{\,4}\times S_4^{\,9}
\]
（推论 \ref{cor:fold-bin-gauge-m6}），并记
\[
A:=S_2^{\,8}\cong(\ZZ_2)^8,\qquad H:=S_3^{\,4}\times S_4^{\,9}.
\]
则 \(Z(H)=1\)、\(Z(\mathcal G^{\mathrm{bin}}_6)=A\)，且 \(A\) 为特征子群。
存在自然同构
\[
\Aut(\mathcal G^{\mathrm{bin}}_6)
\cong
\bigl(\Aut(A)\ltimes\Hom(H,A)\bigr)\rtimes\Aut(H).
\]
进一步，
\[
\Aut(A)\cong \GL(8,2),\qquad
\Hom(H,A)\cong(\ZZ_2)^{104},
\]
并且
\[
\Out(\mathcal G^{\mathrm{bin}}_6)
\cong
\bigl(\GL(8,2)\ltimes(\ZZ_2)^{104}\bigr)\rtimes(S_4\times S_9).
\]
\end{theorem}
\begin{proof}
由推论 \ref{cor:fold-bin-gauge-m6}，\(\mathcal G^{\mathrm{bin}}_6=A\times H\)。
因 \(Z(S_3)=Z(S_4)=1\)，得 \(Z(H)=1\)，故 \(Z(\mathcal G^{\mathrm{bin}}_6)=A\)；中心在自同构下保持，故 \(A\) 为特征子群。

对 \(G=A\times H\)（\(A\) 阿贝尔且中心，\(Z(H)=1\)）的标准分解表明任意自同构可唯一写成
\[
(a,h)\longmapsto \bigl(\alpha(a)\beta(h),\ \delta(h)\bigr),
\]
其中 \(\alpha\in\Aut(A)\)、\(\beta\in\Hom(H,A)\)、\(\delta\in\Aut(H)\)，从而得到第一式。

由 \(A\cong(\ZZ_2)^8\) 得 \(\Aut(A)\cong\GL(8,2)\)。
又
\[
\Hom(H,A)\cong\Hom(H^{\mathrm{ab}},A),\qquad
H^{\mathrm{ab}}\cong(\ZZ_2)^{4+9}=(\ZZ_2)^{13},
\]
故
\[
\Hom(H,A)\cong\Hom((\ZZ_2)^{13},(\ZZ_2)^8)\cong(\ZZ_2)^{104}.
\]

最后，\(\Inn(\mathcal G^{\mathrm{bin}}_6)=\Inn(H)\)（因 \(A\) 中心）。
而 \(S_3,S_4\) 为 complete 群（中心平凡且外自同构平凡），故
\[
\Out(H)\cong S_4\times S_9
\]
仅来自同构因子的置换。
对前述自同构分解取商即得
\[
\Out(\mathcal G^{\mathrm{bin}}_6)
\cong
\bigl(\GL(8,2)\ltimes(\ZZ_2)^{104}\bigr)\rtimes(S_4\times S_9).
\]
证毕。
\end{proof}

\begin{conjecture}[推前非平凡谱的满对称 Galois 普适性（最大不可整合重述）]\label{conj:terminal-window-pushforward-maximal-galois-universality}
对每个 \(m\ge 5\)，令 \(\chi_{P_m}(\lambda)=(\lambda-1)q_m(\lambda)\)。
猜想 \(q_m\) 在 \(\QQ\) 上不可约，且
\[
\Gal(q_m/\QQ)\cong S_{\deg(q_m)}.
\]
该断言与猜想 \ref{conj:terminal-window-pushforward-full-symmetric-spectrum} 同值；
等价地，除平凡本征值 \(1\) 外，全部特征值构成单一 Galois 轨道，且不存在低阶代数关系诱导的系统性可分解结构。
有限窗口 \(m=5,6,7,8,9\) 的证书见表 \ref{tab:window_pushforward_markov_full_symmetric_galois_audit_m5_m9}。
\end{conjecture}

\begin{corollary}[跨尺度二值手征保持的对称破缺必需性]\label{cor:terminal-cross-scale-binary-chirality-symmetry-breaking}
\leanverified{paper\_terminal\_cross\_scale\_binary\_chirality\_symmetry\_breaking}
设对每个尺度 \(m\) 的 boundary 纤维都要求给出符号协变二值标号，并要求该标号与尺度约化相容。
则在二层纤维上该结构可实现；一旦某尺度的 boundary 纤维进入三层化（在 \(m=7,8\) 的审计窗口中已出现），单纤维层面的符号协变标号即不存在。
因此，任何仍欲维持二值手征读出的跨尺度方案，必须引入打破三层纤维 \(S_3\) 全对称性的额外选择数据（例如选定优先二层子对）；该数据在本文语义下属于协议门/外置语义输入。
\end{corollary}
\begin{proof}
由命题 \ref{prop:terminal-sign-covariant-label-boundary-bifurcation}，符号协变二值标号存在当且仅当纤维层数为 \(2\)。
尺度相容条件不能消去单纤维存在性的必要条件；故当纤维层数为 \(3\) 时，必须先引入额外非对称选择，方可将不可协变结构转化为协变读出。证毕。
\end{proof}
